% page 71 du fac-simile (image p-070 du scan)
% bloc 157.5 x 242.0 mm ; marge gauche 44.5 mm
\pgc{22.860mm}{0.000mm}{
\l{}
\l{}
\l{}
\l{}
\l{}
\l{}
\l{\sou{berlino}. Granda karoso klozita, quar-rota, di qua la parti}
\l{\cel{3}avana e posa esas simetra. - DEFIRS.}
\l{}
\l{\sou{berloko}. Siglili, klefi, mikra juveli, qui pendas de rubando,}
\l{\cel{3}de kateno di posh-horlojo. - DFIR.}
\l{}
\l{\sou{bermo}.\cel{2}Parto horizontala en taluso. - DEFRS.}
\l{}
\l{\sou{bernaklo}.\cel{2}(\sou{zool.)}\cel{1}Genero de krustacei kun pedi apendicoza,}
\l{\cel{3}di qui la konko subtenesas da pedikulo kontraktebla. -}
\l{\cel{3}L.\cel{1}\sou{branta bernicla}. - DEFIS.}
\l{}
\l{\sou{bersar}.\cel{2}(\sou{trans.}) Ociligar dum kelka tempo, dolce (infanteto)}
\l{\cel{3}en lua lito o sur la gremio, por dormigar o kalmigar lu.-F.}
\l{}
\l{\sou{bestio}.\cel{2}Omna animalo, eceptante la homo. - DEFIRS.}
\l{}
\l{\sou{beto.}\cel{2}(\sou{bot.)}\cel{2}Planto herbacea, kun radiko pulpoza, genero di}
\l{\cel{3}la familio \textquotedbl{}kenopodei\textquotedbl{}. L.\cel{1}\sou{beta}. - DEFIS.}
\l{}
\l{\sou{betelo}. (\sou{bot.}) Arbusto sarmentoza de India orientala, de la}
\l{\cel{3}genero \textquotedbl{}pipro\textquotedbl{}. L.\cel{1}\sou{piper betel}. - DEFIRS.}
\l{}
\l{\sou{betono}. Mixuro de mortero hidraulikala e de stoni triturita,}
\l{\cel{3}qua hardeskas en aquo. - DEFIRS.}
\l{}
\l{\sou{betonio}. (\sou{bot.}) Planto herbacea, perena, kun bela flori spike-}
\l{\cel{3}dispozite - genero di la familio \textquotedbl{}labiacei\textquotedbl{}. L.\cel{1}\sou{betonica}}
\l{\cel{3}\sou{officinalis}. - DEFIS.}
\l{}
\l{\sou{betravo}. (\sou{bot.}) Speco di beto kun radiko rotacanta quale la}
\l{\cel{3}rapo, blanka o reda, qua donas nutrivo por la homo e la}
\l{\cel{3}bestii, e di qua onu extraktas sukro.\cel{2}L.\cel{1}\sou{beta rapa.}- ES.}
\l{}
\l{\sou{bezigo}. Karto-ludo. - Nomo, en ta ludo, di la unioneso di la}
\l{\cel{3}karelo-pajo e di la piquo-damo. - DEFIS.}
\l{}
\l{\sou{bezoaro}. Agreguro kalkoloza quan onu renkontras en la intes-}
\l{\cel{3}tino ed en la stomako di ula animali, e quan onu uzis kom}
\l{\cel{3}antidoto, amuleto, e c. - DEFR.}
\l{}
\l{\sou{bezonar}.\cel{2}(\sou{trans.}) Indijar ulo quon postulas la naturo o la}
\l{\cel{3}cirkonstanci o la skopo vizata. - FI.}
\l{}
\l{\sou{bi-}\cel{1}.\cel{2}Du-... (uzata por la termini ciencala :\cel{1}\sou{bikapa,}\cel{1}\sou{biloba,}}
\l{\cel{3}e c.}
\l{}
\l{\sou{biblo}. Asemblajo ek la libri \textquotedbl{}Olda Testamento\textquotedbl{} e \textquotedbl{}Nova Testa-}
\l{\cel{3}mento\textquotedbl{} - che la Kristani ; la \textquotedbl{}Olda Testamento\textquotedbl{}, che la}
\l{\cel{3}Judi. - DEFIRS.}
}
